\documentclass[a4paper, 12pt]{article}
\usepackage[T2A]{fontenc}
\usepackage[utf8]{inputenc}
\usepackage[russian]{babel}
\usepackage{array}
\usepackage{booktabs}

\begin{document}
	
	\section{Эксперименты с таблицами}
	
	% 1. Использование простого примера 
	% Создание простейшей таблицы через окружение tabular с двумя колонками, выровненными по левому краю (ll), и горизонтальной чертой \hline.
	\subsection{Простая таблица}
	Я начал с базового примера окружения \texttt{tabular}:
	\begin{center}
		\begin{tabular}{ll}
			Объект & Категория \\
			\hline
			Текст & Данные \\
		\end{tabular}
	\end{center}
	
	% 2. Различные типы выравнивания (l, c, r) (手册第43页)
	% Выравнивание (l, c, r): Демонстрация позиционирования текста внутри ячеек. Использование спецификаторов l (left), c (center) и r (right) в сочетании с @{\hspace{...}} для наглядного отображения пустого пространства.
	\subsection{Демонстрация выравнивания (l, c, r)}
	Я создал три примера, чтобы увидеть разницу в позиционировании текста:
	
	% Слева (l)
	\textbf{1. Выравнивание по левому краю (l):} \\
	\begin{tabular}{|l@{\hspace{3cm}}|}
		\hline
		Текст слева \\
		\hline
	\end{tabular}
	
	\vspace{0.5cm}
	
	% По центру (c)
	\textbf{2. Выравнивание по центру (c):} \\
	\begin{tabular}{|@{\hspace{1.5cm}}c@{\hspace{1.5cm}}|}
		\hline
		Текст по центру \\
		\hline
	\end{tabular}
	
	\vspace{0.5cm}
	
	% Справа (r)
	\textbf{3. Выравнивание по правому краю (r):} \\
	\begin{tabular}{|@{\hspace{3cm}}r|}
		\hline
		Текст справа \\
		\hline
	\end{tabular}
	
	% 3. Слишком мало элементов 
	% Недостаток элементов: Анализ поведения компилятора при пропуске данных. Если количество символов & меньше объявленных колонок, LaTeX просто оставляет правую часть строки пустой.
	\subsection{Слишком мало элементов}
	Если в строке меньше элементов, чем объявлено в преамбуле (например, \texttt{ccc}), LaTeX просто оставляет пустые ячейки справа.
	\begin{center}
		\begin{tabular}{|c|c|c|}
			\hline
			Первый & Второй & \\ % Третий элемент пропущен
			\hline
		\end{tabular}
	\end{center}
	
	%4. Слишком много элементов 
	%\subsection{Слишком много элементов}
    	%Если в строке больше элементов (лишние символы \texttt{\&}), LaTeX выдаст ошибку \texttt{Extra alignment tab has been changed to \cr}. Это происходит %потому, что количество разделителей превышает число объявленных колонок.	
	
	%5. Команда \multicolumn 
	%Объединение ячеек: Использование команды \multicolumn{n}{pos}{text} для слияния нескольких колонок в одну с изменением типа выравнивания для конкретной строки
	\subsection{Команда \textbackslash\textbackslash multicolumn}
	Я использовал эту команду, чтобы объединить несколько колонок в одну:
	\begin{center}
		\begin{tabular}{lll}
			\toprule
			Колонка 1 & Колонка 2 & Колонка 3 \\
			\midrule
			\multicolumn{2}{c}{Объединенные ячейки} & Одиночная \\
			\bottomrule
		\end{tabular}
	\end{center}
	
\end{document}